\subsection{An affine primitive for act-cone products}

\begin{theorem}[A common fixed row for nonconstant column actors]
\label{C:act-common-fixed-row-cups}
Let $M$ have rectangular-band minimal ideal $I\times\{0,1\}$.
Write $\sigma_m:I\to I$ and
$\tau_m:\{0,1\}\to\{0,1\}$ for its row and column actions.
Suppose there is a row $i_0$ fixed by every $\sigma_m$ for which
$\tau_m$ is nonconstant.  Then for every nonempty finite right
$M$-set $S$, all products
\[
 H^2(B\mathcal C(M,S);\mathbb Z)\otimes
 H^2(B\mathcal C(M,S);\mathbb Z)
       \longrightarrow H^4(B\mathcal C(M,S);\mathbb Z)
\]
are zero.  In particular no such act cone is a projective plane.
There is no bound on the number of rows, states, or outside elements,
and no finiteness hypothesis on the classifying space.
\end{theorem}
\begin{proof}
By Lemma~\ref{C:constant-action-quotient-H2}, every integral
degree-two class pulls back from the constant-action quotient
$\bar S$.  Compress that quotient to a monoid $N$ with minimal
band $I\times\bar S$, as in
Theorem~\ref{C:constant-minimal-act-compression}.
Let $s_j\in\bar S$ be the constant value of an original band
element in column $j$.  This is independent of its row, and
\[
 s_jm=s_{\tau_m(j)}.
\]
The two values are allowed to coincide.  Use the base row $i_0$
and base column $s_0$ throughout.

Put $V=\widetilde{\mathbb ZI}$ and
$W=\widetilde{\mathbb Z\bar S}$.
By Lemma~\ref{lem:B-band-Tor}, integral degree-two classes on
$BN$ are precisely integral balanced bilinear forms
$L:W\times V\to\mathbb Z$.
Here $H_1(BN)=0$, so the universal-coefficient identification
has no additional degree-two extension term.
For an actor $f$ of $N$ define
\[
 a_f=[\sigma_f(i_0)]-[i_0]\in V,\qquad
 b^L_f(v)=L([\tau_f(s_0)]-[s_0],v).
\]
Let $A_f$ be its row augmentation operator.  The affine identities
and balance give
\[
 a_{fg}=a_f+A_fa_g,\qquad
 b^L_{fg}=b^L_g+b^L_fA_g.
\]
Consequently $u_L(f,g)=b^L_fa_g$ is the usual integral
two-cocycle representing $L$.

We make its transfer to the act category literal.
Put $e=(i_0,s_0)$ in the band of $N$, and take the full
two-object Karoubi subcategory on $1,e$.
The compression functor sends the monoid object to $1$, the
sink to $e$, an old actor $m$ to its image in $N$, and a cross
state $s$ to $(i_0,s)$.  Extend $u_L$ to this subcategory by the
same formula on the underlying $N$ labels.
Both object identities $1,e$ have $a=b^L=0$, so this is a
normalized category cocycle; its coboundary is zero by the
same affine identities.  Its restriction to the object $1$
is the original monoid cocycle.  The inclusion of this object
is a classifying-space equivalence, so the extension represents
the unique corresponding cohomology class.  Pulling back along
the honest compression functor gives the desired representative
on $\mathcal C(M,\bar S)$.
For a cross arrow its row vector $a$ is zero and its column
functional is $L([s]-[s_0],-)$.  All other labels are old actors.

Define the column character and translation by
\[
 \chi(m)=\tau_m(1)-\tau_m(0),\qquad
 \beta(m)=\tau_m(0).
\]
They obey
\[
 \beta(mn)=\beta(n)+\chi(n)\beta(m).
\]
For each balanced form $L$, put
\[
 \ell_L(v)=L([s_1]-[s_0],v),\qquad
 \alpha_L(m)=\ell_L(a_m).
\]
The equivariant map from the original column augmentation line
to $W$ gives $\ell_LA_m=\chi(m)\ell_L$.  Therefore
\[
 \alpha_L(mn)=\alpha_L(m)+\chi(m)\alpha_L(n),
 \qquad
 b^L_m=\beta(m)\ell_L,
 \qquad
 u_L(m,n)=\beta(m)\alpha_L(n)
\]
on old monoid arrows.
The common-fixed-row hypothesis is exactly
\[
 \chi(m)a_m=0\qquad(m\in M).
\]

Use algebra-ordered normalized chains, so their adjacent
composition is multiplication from left to right.  A chain in
an act category either consists entirely of old monoid arrows
or has one cross arrow in its first position; all subsequent
arrows belong to $M$.  For two balanced forms $L,D$ define an
integral three-cochain to be zero on every word containing an
identity, and on the nondegenerate typed domain define it by
\[
 P(f,g,h)=u_L(f,g)\,\beta(g)\,\alpha_D(h).
\]
This prescription is normalized; on a nondegenerate word the
arguments of $\beta$ and $\alpha_D$ are always old actors.
Expand its five-term coboundary on
$(f,g,h,k)$ and use the two-cocycle identity for $u_L$:
\begin{align*}
 (\delta P)(f,g,h,k)
 &=u_L(f,g)\beta(h)\alpha_D(k)\\
 &\quad+\chi(h)\beta(g)\alpha_D(k)
       \bigl(u_L(f,gh)-u_L(f,g)\bigr)\\
 &=u_L(f,g)u_D(h,k)
   +\chi(h)\beta(g)\alpha_D(k)b^L_f A_g a_h\\
 &=u_L(f,g)u_D(h,k).
\end{align*}
The last equality uses $\chi(h)a_h=0$.
Thus $\delta P=u_L\smile u_D$ as an equality of integral
normalized category cochains.  Every degree-two class is
represented this way after the quotient pullback, proving the
assertion about all mixed products.
\end{proof}

\begin{corollary}[All dual-number three-row act cones]
\label{C:dual-row-all-act-cups}
Suppose $M$ has a $3\times2$ rectangular-band minimal ideal and
the image of its rational row augmentation representation is
the nonsemisimple dual-number algebra.
For every finite right $M$-set $S$, all integral degree-two
products on $B\mathcal C(M,S)$ are zero.
In particular no such cone realizes $\mathbb CP^2$, regardless
of the sizes of $M$ and $S$.
\end{corollary}
\begin{proof}
Normalize the row labels as in
Proposition~\ref{B:three-row-dual-cup-radical}.
The actual row transformations are
\[
 (0,1,2),\quad(0,0,1),\quad(1,1,0),\quad
 (0,0,0),\quad(1,1,1),\quad(2,2,2).
\]
Their augmentation operators are $I,N,-N,0,0,0$, with
$N=\left(\begin{smallmatrix}0&1\\0&0\end{smallmatrix}\right)$.
Write $\psi$ for their common diagonal character; it is one
only for the identity row map and zero otherwise.
The only nonzero one-dimensional quotient of the rational row
module has character $\psi$, and every such quotient kills
$\operatorname{im}N$.

Retain the notation of the preceding proof for an arbitrary
integral balanced form $L$ after constant-action compression.
Its old-column covector $\ell_L$ satisfies
$\ell_LA_m=\chi(m)\ell_L$.
If the original column character $\chi$ differs from $\psi$,
every such covector is zero, since no quotient of that character
exists over $\mathbb Q$.  Hence every transferred cocycle
$u_L(m,n)$ is literally zero on pairs of old monoid arrows.
On every four-chain the second factor of a degree-two cup
product is such a pair; all products are therefore already
zero as cochains.

If $\chi=\psi$, each nonconstant column actor has $\chi=1$,
and its row transformation is the identity.  Every row is a
common fixed row, and
Theorem~\ref{C:act-common-fixed-row-cups} applies.
Both arguments use integral cochains and the integral
degree-two quotient surjection; the rational representation
classification introduces no denominator into the primitive.
\end{proof}

\begin{remark}[A completed finite control and a superseded search]
\label{C:dual-row-act-control-scope}
Before the preceding primitive was found, a bounded upper-model
search was completed for one explicit six-dimensional contracted
basis: a square-zero loop at the first vertex of a zero-composite
three-vertex chain.  Its $72$ literal $3\times2$ affine lifts and
$144$ trace/rank cases gave $593536$ eight-state acts after the
trace constraints.  Of these, $275040$ had the required exact
Tor profile, and $90240$ retained a possible rank-two balanced
class.  The $10236$ state-orbit representatives all had zero
mod-two square, with actual stored primitives checked on
$258858204$ four-simplices.
The code and outputs have prefix \path{a_dual_loop_chain}.
This finite control is not used in the all-order proof.

The broader restricted order-twelve dual-row target stream
\path{code/outputs/a_twelve_dual_targets.tsv} was intentionally
interrupted when the all-order theorem superseded it.  It is
retained as an incomplete research prefix, not a complete census
and not a premise of any exclusion.
\end{remark}
